\section{结果与讨论}
\label{sec:results-discussion}

本章基于正式结果包对 CVF-AE 的性能进行分析。讨论顺序为：首先给出 11 个算法在三类低空路径规划场景中的主对比结果，然后分析统计显著性、消融结果、评价开销、参数稳健性和轨迹合理性。由于本文任务是强约束低空 UAV 路径规划，单一标量适应度并不能完全代表工程可用性。因此，下文将适应度、可行率、最终违反量、轨迹复核指标、评价次数和 CVF 触发规模联合解释。

\subsection{主对比结果}

\begin{table*}[tbp]
\centering
\caption{主对比实验结果}
\label{tab:main-comparison}
\scriptsize
\setlength{\tabcolsep}{3pt}
\resizebox{\textwidth}{!}{%
\begin{tabular}{lccccccccccc}
\toprule
Algorithm & S1 $J$ & S1 FR & S1 R & S2 $J$ & S2 FR & S2 R & S3 $J$ & S3 FR & S3 R & Avg. R & Eng. flag \\
\midrule
AE & 656.1972 $\pm$ 325.9113 & 0.60 & 7 & 1627.2324 $\pm$ 330.2622 & 0.00 & 7 & 1680.4468 $\pm$ 208.6819 & 0.00 & 7 & 7.00 & 0.22 \\
PSO & 423.3491 $\pm$ 216.2313 & 0.90 & 6 & 659.5587 $\pm$ 303.2289 & 0.43 & 6 & 1303.4640 $\pm$ 388.2979 & 0.03 & 6 & 6.00 & 0.49 \\
GWO & 335.7826 $\pm$ 32.9433 & 0.97 & 4 & 341.1693 $\pm$ 41.4608 & 1.00 & 4 & 673.2786 $\pm$ 376.6002 & 0.53 & 4 & 4.00 & 0.93 \\
HHO & 361.7454 $\pm$ 71.7330 & 0.97 & 5 & 357.9223 $\pm$ 149.3598 & 0.97 & 5 & 929.5468 $\pm$ 437.4162 & 0.23 & 5 & 5.00 & 0.78 \\
ERIME & 318.5440 $\pm$ 15.4114 & 1.00 & 3 & 320.8775 $\pm$ 24.5106 & 1.00 & 3 & 356.4928 $\pm$ 24.6010 & 1.00 & 2 & 2.67 & 0.79 \\
MSCSO & 307.8580 $\pm$ 14.8041 & 1.00 & 2 & 295.0138 $\pm$ 19.1301 & 1.00 & 1 & 449.9550 $\pm$ 198.8820 & 0.83 & 3 & 2.00 & 0.64 \\
CVF-AE & 298.5286 $\pm$ 2.6211 & 1.00 & 1 & 301.1934 $\pm$ 2.8253 & 1.00 & 2 & 350.7658 $\pm$ 23.3272 & 1.00 & 1 & 1.33 & 0.10 \\
\bottomrule
\end{tabular}%
}
\vspace{1mm}
\parbox{\textwidth}{\footnotesize 注：FR 为可行率，R 为主对比算法按平均适应度升序得到的场景排名，Avg. R 为三场景平均排名。Eng. flag 为三场景共 90 次运行中被任一工程轨迹复核规则标记的比例；该固定阈值复核仅作 sanity check，不作为安全认证指标。}
\end{table*}


表 \ref{tab:main-comparison} 给出了扩展主对比结果。CVF-AE 在三个场景中的可行率均为 $1.00$，平均场景排名为 $2.00$，在 11 个算法中获得最佳平均场景排名。该结果说明，CVF-AE 的主要优势不是在每个场景都取得最低标量适应度，而是在强约束环境下同时保持可行性、较低目标值和稳定轨迹质量。

\begin{figure}[tbp]
\centering
\includegraphics[width=\linewidth]{cvf_ae_boxplots_strong_algorithms.png}
\caption{强竞争算法的最终适应度分布}
\label{fig:boxplots-strong}
\end{figure}

图 \ref{fig:boxplots-strong} 展示了 CVF-AE 与强竞争算法的逐次运行分布。CVF-AE 在 Scene 1 和 Scene 3 中保持较低且较稳定的适应度分布；Scene 2 中 CPO 和 MSCSO 的箱体位置更低，对应表 \ref{tab:significance} 中的显著劣势，需要结合轨迹复核进一步解释。

在 Scene 1 中，CVF-AE 的平均适应度为 $298.1379$，略低于 CPO 的 $298.3594$，二者均达到 $1.00$ 的可行率。该场景相对宽松，强约束引导的优势主要体现在保持低风险、低绕行和稳定可行性，而不是拉开大幅标量差距。AE、WOA、HHO 和 PSO 在该场景中均出现较明显的可行率下降或适应度劣化，说明单纯基础进化搜索或无约束场机制难以稳定处理三维障碍、禁飞区和高度限制的组合。

在 Scene 2 中，CPO 和 MSCSO 的标量适应度分别为 $278.1196$ 和 $291.1680$，低于 CVF-AE 的 $301.2018$。因此，不能将 CVF-AE 写成该场景的标量最优方法。然而，轨迹复核显示 CPO 在该场景的高空复核标记率为 $1.00$，MSCSO 的高空复核标记率为 $0.50$，而 CVF-AE 的高空和边界复核标记率均为 $0$。这说明，较低的 scalar fitness 可能来自更激进的高度或通道选择，未必代表更符合低空走廊任务的路径。CVF-AE 在该场景提供的是更保守且稳定的可行路径，而不是单纯追求最低目标值。

在 Scene 3 中，GDSAO 的平均适应度为 $346.8474$，略低于 CVF-AE 的 $349.2270$；但 GDSAO 的边界贴靠复核标记率达到 $0.9667$，ERIME 也达到 $0.9667$，DBO 为 $0.9000$。相比之下，CVF-AE 的 Scene 3 边界贴靠复核标记率为 $0.1667$，且可行率保持 $1.00$。因此，Scene 3 的结果进一步表明，强约束低空路径规划不能只依据最终适应度排序，还需要检查路径是否过度贴近空间边界、是否绕开了实际低空走廊以及是否保持可行。

图 \ref{fig:convergence} 给出了三类场景的收敛行为。CVF-AE 在 Scene 1 和 Scene 3 中能够较快进入稳定低适应度区间；Scene 2 中 CPO 的最终均值更低，进一步说明本文结论应聚焦综合工程表现，而不是单一收敛终值。

\begin{figure}[tbp]
\centering
\begin{subfigure}{0.96\linewidth}
    \centering
    \includegraphics[width=\linewidth]{scene1_convergence_mean_std.png}
    \caption{Scene 1}
\end{subfigure}

\begin{subfigure}{0.96\linewidth}
    \centering
    \includegraphics[width=\linewidth]{scene2_convergence_mean_std.png}
    \caption{Scene 2}
\end{subfigure}

\begin{subfigure}{0.96\linewidth}
    \centering
    \includegraphics[width=\linewidth]{scene3_convergence_mean_std.png}
    \caption{Scene 3}
\end{subfigure}
\caption{三类测试场景的平均收敛曲线}
\label{fig:convergence}
\end{figure}

图 \ref{fig:extended-paths-top} 和图 \ref{fig:extended-paths-3d} 分别从俯视视角和三维视角展示轨迹层面的差异。三维视图用于观察高度变化和障碍物空间关系，俯视轨迹更容易暴露边界贴靠、过度绕行和局部通道选择问题。与只看适应度均值相比，两类视图共同说明了本文在 Scene 2 和 Scene 3 中对 CPO、GDSAO、ERIME 与 MSCSO 保持谨慎解释的原因。

\subsection{统计显著性分析}

\begin{table}[tbp]
\centering
\caption{CVF-AE 与基线算法的逐场景显著性比较}
\label{tab:significance}
\scriptsize
\setlength{\tabcolsep}{2.8pt}
\begin{tabular}{@{}llccc@{}}
\toprule
Baseline & Scene & $p_{\mathrm{Holm}}$ & $\delta$ & R \\
\midrule
AE & 1 & 3.0e-10 & +1.00 & + \\
AE & 2 & 3.0e-10 & +1.00 & + \\
AE & 3 & 3.0e-10 & +1.00 & + \\
PSO & 1 & 2.0e-07 & +0.81 & + \\
PSO & 2 & 5.5e-04 & +0.59 & + \\
PSO & 3 & 3.0e-10 & +1.00 & + \\
GWO & 1 & 3.0e-10 & +1.00 & + \\
GWO & 2 & 0.001 & +0.55 & + \\
GWO & 3 & 2.1e-05 & +0.70 & + \\
WOA & 1 & 3.0e-10 & +1.00 & + \\
WOA & 2 & 3.0e-10 & +1.00 & + \\
WOA & 3 & 3.0e-10 & +1.00 & + \\
HHO & 1 & 1.0e-09 & +0.95 & + \\
HHO & 2 & 0.001 & +0.55 & + \\
HHO & 3 & 5.9e-08 & +0.87 & + \\
DBO & 1 & 3.0e-10 & +1.00 & + \\
DBO & 2 & 1.5e-06 & +0.78 & + \\
DBO & 3 & 5.1e-05 & +0.66 & + \\
CPO & 1 & 0.773 & +0.04 & = \\
CPO & 2 & 3.0e-10 & -1.00 & - \\
CPO & 3 & 0.934 & +0.14 & = \\
GDSAO & 1 & 5.5e-10 & +0.97 & + \\
GDSAO & 2 & 0.620 & -0.08 & = \\
GDSAO & 3 & 0.934 & +0.15 & = \\
ERIME & 1 & 4.0e-10 & +0.98 & + \\
ERIME & 2 & 0.595 & +0.16 & = \\
ERIME & 3 & 0.097 & +0.34 & = \\
MSCSO & 1 & 4.8e-07 & +0.78 & + \\
MSCSO & 2 & 0.002 & -0.51 & - \\
MSCSO & 3 & 0.934 & +0.13 & = \\
\bottomrule
\end{tabular}
\end{table}


表 \ref{tab:significance} 给出了主对比中 CVF-AE 与每个基线的逐场景 Holm 校正 $p$ 值和 Cliff's delta。三个场景的 Kruskal--Wallis 总体检验均显示算法间分布存在显著差异；进一步的成对比较采用双侧 Wilcoxon rank-sum 检验，并在每个场景内使用 Holm 校正。主对比中，CVF-AE 对 10 个基线算法的逐场景比较共得到 $21$ 胜、$7$ 平和 $2$ 负。两项显著劣势均出现在 Scene 2，分别对应 CPO 和 MSCSO 的标量适应度优势。表中 $R$ 表示 Holm 校正后的胜、平、负关系；$\delta$ 按最小化问题定义，正值表示 CVF-AE 的最终适应度分布更优，负值表示对应基线更优。

该统计结果支持两个结论。第一，CVF-AE 在多数场景和多数基线上具有显著或至少不劣的适应度表现，尤其相对于 AE、PSO、GWO、WOA、HHO 和 DBO 等基线表现稳定。第二，统计检验本身只作用于最终最佳适应度分布，不能替代轨迹合理性判断。Scene 2 中 CPO 和 MSCSO 的标量优势是真实存在的，但结合高空复核标记和 Scene 3 的可行率/边界贴靠风险后，本文更适合将 CVF-AE 的结论表述为综合均衡优势，而不是逐场景标量最优。

\begin{figure*}[tbp]
\centering
\begin{subfigure}{0.32\linewidth}
    \centering
    \includegraphics[width=\linewidth]{extended_scene1_representative_top.png}
    \caption{Scene 1}
\end{subfigure}
\hfill
\begin{subfigure}{0.32\linewidth}
    \centering
    \includegraphics[width=\linewidth]{extended_scene2_representative_top.png}
    \caption{Scene 2}
\end{subfigure}
\hfill
\begin{subfigure}{0.32\linewidth}
    \centering
    \includegraphics[width=\linewidth]{extended_scene3_representative_top.png}
    \caption{Scene 3}
\end{subfigure}
\caption{三类测试场景的代表性路径俯视对比}
\label{fig:extended-paths-top}
\end{figure*}

\begin{figure*}[tbp]
\centering
\begin{subfigure}{0.32\linewidth}
    \centering
    \includegraphics[width=\linewidth]{extended_scene1_representative_3d.png}
    \caption{Scene 1}
\end{subfigure}
\hfill
\begin{subfigure}{0.32\linewidth}
    \centering
    \includegraphics[width=\linewidth]{extended_scene2_representative_3d.png}
    \caption{Scene 2}
\end{subfigure}
\hfill
\begin{subfigure}{0.32\linewidth}
    \centering
    \includegraphics[width=\linewidth]{extended_scene3_representative_3d.png}
    \caption{Scene 3}
\end{subfigure}
\caption{三类测试场景的代表性路径三维对比}
\label{fig:extended-paths-3d}
\end{figure*}

\subsection{消融实验}

\begin{table*}[tbp]
\centering
\caption{CVF-AE 关键模块消融实验结果}
\label{tab:ablation}
\setlength{\tabcolsep}{3pt}
\tiny
\resizebox{\textwidth}{!}{%
\begin{tabular}{clcccccc}
\toprule
Scene & Variant & $J$ & FR & R & Avg. R & nEvals & CVF trig./succ. \\
\midrule
1 & CVF-AE & 297.8318 $\pm$ 3.1119 & 1.00 & 1 & 1.00 & 9566.4 & 21.1/19.2 \\
1 & CVF-AE-w/o-CVF & 306.1985 $\pm$ 3.4917 & 1.00 & 3 & 3.00 & 9471.8 & 0.0/0.0 \\
1 & CVF-AE-w/o-StateAdaptiveCVF & 298.3037 $\pm$ 3.2161 & 1.00 & 2 & 2.33 & 9566.0 & 20.6/19.7 \\
1 & CVF-AE-w/o-SparsePreservation & 306.3793 $\pm$ 4.9973 & 1.00 & 4 & 3.67 & 9140.1 & 110.1/90.7 \\
1 & CVF-AE-w/o-Init & 316.9964 $\pm$ 11.8481 & 1.00 & 5 & 5.00 & 9587.6 & 140.2/93.1 \\
1 & Base-AE & 654.0475 $\pm$ 328.7142 & 0.63 & 6 & 6.00 & 9030.0 & 0.0/0.0 \\
2 & CVF-AE & 300.6612 $\pm$ 3.2775 & 1.00 & 1 & 1.00 & 9551.6 & 10.6/7.9 \\
2 & CVF-AE-w/o-CVF & 303.7697 $\pm$ 4.9565 & 1.00 & 4 & 3.00 & 9432.5 & 0.0/0.0 \\
2 & CVF-AE-w/o-StateAdaptiveCVF & 301.8164 $\pm$ 2.9332 & 1.00 & 2 & 2.33 & 9560.9 & 20.4/14.6 \\
2 & CVF-AE-w/o-SparsePreservation & 303.3114 $\pm$ 6.0718 & 1.00 & 3 & 3.67 & 9107.8 & 77.8/51.5 \\
2 & CVF-AE-w/o-Init & 313.2191 $\pm$ 18.3412 & 1.00 & 5 & 5.00 & 9592.4 & 153.7/107.4 \\
2 & Base-AE & 1597.4829 $\pm$ 342.4791 & 0.00 & 6 & 6.00 & 9030.0 & 0.0/0.0 \\
3 & CVF-AE & 348.5856 $\pm$ 21.8153 & 1.00 & 1 & 1.00 & 9367.1 & 86.8/59.5 \\
3 & CVF-AE-w/o-CVF & 355.8347 $\pm$ 21.3086 & 1.00 & 2 & 3.00 & 9258.5 & 0.0/0.0 \\
3 & CVF-AE-w/o-StateAdaptiveCVF & 361.3123 $\pm$ 35.0403 & 1.00 & 3 & 2.33 & 9321.8 & 72.5/46.3 \\
3 & CVF-AE-w/o-SparsePreservation & 372.2255 $\pm$ 19.8216 & 1.00 & 4 & 3.67 & 9118.7 & 88.7/56.6 \\
3 & CVF-AE-w/o-Init & 556.7515 $\pm$ 348.8264 & 0.70 & 5 & 5.00 & 9802.7 & 417.9/336.5 \\
3 & Base-AE & 1643.7320 $\pm$ 199.4027 & 0.00 & 6 & 6.00 & 9030.0 & 0.0/0.0 \\
\bottomrule
\end{tabular}%
}
\vspace{1mm}
\parbox{\textwidth}{\footnotesize 注：FR 为可行率，R 为场景内排名，Avg. R 为跨场景平均排名。CVF trig./succ. 表示平均触发的 CVF 候选数 / 平均被局部接收的 CVF 候选数。}
\end{table*}


表 \ref{tab:ablation} 展示了正式消融实验结果。完整 CVF-AE 在三个场景中的消融平均排名为 $1.00$，并在所有场景保持 $1.00$ 的可行率。与 Base-AE 相比，完整方法在 Scene 2 和 Scene 3 中将可行率从 $0$ 提升到 $1.00$，说明仅依靠基础 AE 搜索无法稳定穿越强约束低空场景，而约束感知初始化、CVF 候选和稀疏保持机制共同提高了可行路径形成能力。

图 \ref{fig:ablation-scene2-top} 和图 \ref{fig:ablation-scene3-3d} 展示了完整方法与消融版本的路径形态差异。Scene 2 的俯视图说明完整方法更稳定地沿低风险可行走廊推进；去除关键模块后，路径更容易出现绕行、靠近障碍或需要更多后处理修复。Scene 3 的三维图进一步补充高度维度证据：由于高度上界更低、障碍和禁飞区更密集，去除初始化或 CVF 相关机制后，算法更容易依赖较差的空间绕行或停留在低质量可行区域。两图与表 \ref{tab:ablation} 中 Scene 3 的适应度和可行率变化相互印证。

\begin{figure}[tbp]
\centering
\includegraphics[width=0.78\linewidth]{ablation_scene2_representative_top.png}
\caption{Scene 2 关键模块消融版本的代表性路径俯视对比}
\label{fig:ablation-scene2-top}
\end{figure}

\begin{figure}[tbp]
\centering
\includegraphics[width=0.78\linewidth]{ablation_scene3_representative_3d.png}
\caption{Scene 3 关键模块消融版本的代表性路径三维对比}
\label{fig:ablation-paths}
\label{fig:ablation-scene3-3d}
\end{figure}

去除 CVF 场后，算法在三个场景仍能保持可行，但适应度均有所上升：Scene 1 从 $297.8318$ 上升到 $306.1985$，Scene 2 从 $300.6612$ 上升到 $303.7697$，Scene 3 从 $348.5856$ 上升到 $355.8347$。这说明在已有初始化和 repair 支撑下，CVF 不是唯一可行性来源，但能够在可行区域附近提供额外的方向引导，使路径质量进一步改善。

去除约束感知初始化的影响在 Scene 3 最明显，可行率从 $1.00$ 降至 $0.70$，平均适应度升至 $556.7515$，并伴随更多 CVF 触发。这表明初始化并不是简单加速技巧，而是强约束走廊中避免搜索长期滞留不可行区域的重要组成部分。去除稀疏保持机制后，三个场景的适应度均弱于完整方法，说明局部修复与可行性保持对路径稳定性仍有贡献。

状态自适应 CVF 的作用更细致。消融均值显示，去除状态自适应 CVF 后并非每个场景都明显劣化，这意味着状态调度不宜被解释为单独决定性能的模块。更稳妥的解释是：状态自适应、CVF 场、初始化和稀疏保持在不同约束阶段形成互补，完整机制在总体上保持最稳定的综合表现。

\subsection{评价开销与参数稳健性}

\begin{table}[tbp]
\centering
\caption{CVF-AE 评价开销与稀疏触发统计}
\label{tab:runtime}
\tiny
\resizebox{\columnwidth}{!}{%
\begin{tabular}{@{}lcccccc@{}}
\toprule
Group & Sc. & nEvals & Trig. & Succ. & Rate & Ratio \\
\midrule
Main & 1 & 9569.5 & 25.0 & 22.4 & 0.28\% & -- \\
Main & 2 & 9552.7 & 12.4 & 9.5 & 0.14\% & -- \\
Main & 3 & 9362.5 & 83.8 & 56.4 & 0.93\% & -- \\
\midrule
Abl. & 1 & 9566.4 & 21.1 & 19.2 & 0.22\% & 1.010 \\
Abl. & 2 & 9551.6 & 10.6 & 7.9 & 0.11\% & 1.013 \\
Abl. & 3 & 9367.1 & 86.8 & 59.5 & 0.93\% & 1.012 \\
\bottomrule
\end{tabular}%
}
\vspace{1mm}
\parbox{\columnwidth}{\footnotesize 注：Trig. 和 Succ. 分别表示平均触发的 CVF 候选数和平均被接收的 CVF 候选数；Ratio 为完整 CVF-AE 相对于同批次 w/o-CVF 消融版本的评价次数比。}
\end{table}


表 \ref{tab:runtime} 给出了 CVF-AE 的评价开销摘要。主对比中，CVF-AE 在 Scene 1、Scene 2 和 Scene 3 的平均评价次数分别为 $9568.9$、$9554.2$ 和 $9355.2$，与 $N=30$、$T=300$ 的基础评价规模接近。CVF 触发次数在 Scene 1 和 Scene 2 分别为 $23.9$ 和 $13.1$，对应触发率低于 $0.3\%$；在更强约束的 Scene 3 中，触发次数增至 $88.6$，触发率仍低于 $1\%$。该趋势与方法设计一致：当约束更密集、可行走廊更窄时，稀疏 CVF 分支会更频繁地被触发，但不会退化为全种群强制更新。

同批次消融对比显示，完整 CVF-AE 相对于去除 CVF 的版本，在三个场景中的评价次数比值约为 $1.010$、$1.013$ 和 $1.012$。这支持“额外评价次数受控”的结论。本文不将绝对 wall-clock time 作为主要效率证据，因为 MATLAB 运行时间会受到 repair 次数、路径可行形成速度、会话状态和批处理负载影响。更合理的结论是：CVF-AE 的额外评价开销约为 $1\%$ 量级，低开销主要来自稀疏触发和局部候选接收，而不是来自更短的绝对运行时间。

参数稳健性检查进一步作为补充说明使用，而不作为独立性能证据。候选触发配额、CVF 步长强度和保守状态切换阈值的小范围扰动显示：更高触发配额可在 Scene 2 和 Scene 3 中降低部分适应度均值，但会增加触发次数和额外评价规模；步长强度的较优档位随场景变化；状态切换阈值没有在两个检查场景上同时稳定占优的非默认设置。因此，本文保留默认参数，并将结论限定为默认配置具有较好的折中稳健性，而不是依赖单一场景调参获得优势。

\subsection{综合讨论}

综合主对比、显著性、消融、开销和参数稳健性分析，CVF-AE 的价值主要体现在三个方面。第一，在所测试的三类合成低空场景中，它保持 $1.00$ 的可行率，并在扩展主对比中获得最佳平均排名，说明约束可行性场与可行性优先准则能够稳定服务于本文设定的强约束路径规划任务。第二，它在 Scene 2 和 Scene 3 这类更强约束场景中提供了比多数基线更均衡的路径，不以高空绕行或边界贴靠作为主要代价。第三，稀疏触发使 CVF 额外评价次数保持在较低水平，避免了每代全种群构造可行性场所带来的计算膨胀。

本文的结论也有明确边界。CVF-AE 不是每个场景的最低标量适应度算法；Scene 2 中 CPO 和 MSCSO 的适应度显著更低，Scene 3 中 GDSAO 的均值也略低于 CVF-AE。本文将这些结果解释为低空 UAV 路径规划中的多指标折中：最终适应度、可行率、轨迹复核风险和评价开销必须共同考虑。基于这一口径，CVF-AE 在强约束低空路径规划任务中表现为更均衡、更稳定的工程型优化方法。
